\documentclass[11pt]{article}
\usepackage[margin=1in]{geometry}
\usepackage{amsmath,amssymb,amsthm}
\usepackage{booktabs}
\usepackage{hyperref}
\usepackage{listings}
\usepackage{xcolor}
\usepackage{siunitx}

\title{Independent Replication --- Anton, Cohen \& Quer-Sardanyons (2017/2020):\\
A Fully Discrete Stochastic Exponential Integrator (SEXP) for the\\
1D Stochastic Heat Equation}
\author{PDE-100 Replication Wave / Ollie \\ \small (independent replication report)}
\date{\today}

\begin{document}
\maketitle

\begin{abstract}
We independently reproduce the fully-discrete stochastic exponential integrator (SEXP) of
Anton, Cohen \& Quer-Sardanyons (arXiv:1711.08340, IMA J.\ Numer.\ Anal.)\ for the 1D
stochastic heat equation with multiplicative space--time white noise. Using a from-scratch
DST-diagonalized implementation validated to machine precision on the deterministic case,
we confirm (i) CFL-freeness (stability across
$\Delta t \in [2^{-1},2^{-16}]$ at $M=512$), (ii) empirical strong temporal convergence
order $\approx 1/2$ (measured RMS slope \textbf{0.558} over 500 samples at the paper's exact
parameters, matching Fig.~1), and (iii) pathwise (almost-sure) convergence with per-path
order $\approx 1/2$. Three independent free-endpoint LLM judges (\texttt{argo:gpt-5.2},
\texttt{argo:gemini-2.5-pro}, \texttt{argo:gpt-4.1}) return \textsc{Replicated}.
\textbf{Verdict: REPLICATED.}
\end{abstract}

\section{Paper summary}

The paper studies an explicit fully-discrete scheme for the 1D stochastic heat equation
with multiplicative space--time white noise:
\begin{equation}
  \frac{\partial u}{\partial t} \;=\; \frac{\partial^2 u}{\partial x^2}
    \;+\; f(t,x,u) \;+\; \sigma(t,x,u)\,\frac{\partial^2 W}{\partial t\,\partial x},
    \qquad (t,x)\in(0,\infty)\times(0,1),
\end{equation}
$u(t,0)=u(t,1)=0$, $u(0,\cdot)=u_0$, $W$ a Brownian sheet. The space discretization uses
the standard second-difference operator on the interior grid $x_m=m/M$
($m=1,\dots,M-1$, $\Delta x = 1/M$), giving the $(M-1)$-dimensional SDE system
\[
  du^M \;=\; M^2 D\,u^M\,dt \;+\; f\,dt \;+\; \sqrt{M}\,\sigma\,dW^M,
\]
with $D = \mathrm{tridiag}(1,-2,1)$ and $A := M^2 D$. Eigenvalues
$\lambda_j = -4M^2\sin^2(j\pi/2M)$ with sine eigenvectors $\varphi_j(x)=\sqrt{2}\sin(j\pi x)$.
Time is discretized with the \emph{stochastic exponential integrator} (Eq.\ (15)):
\begin{equation}
  U^{n+1} \;=\; e^{A\Delta t}\,U^n \;+\; F(U^n)\,\Delta t \;+\; \Sigma(U^n)\,\Delta W^n,
  \label{eq:sexp}
\end{equation}
where $F(U)_m=f(U_m)$, $\Sigma(U)$ is diagonal with entries $\sqrt{M}\sigma(U_m)$, and
$\Delta W^n$ are $(M-1)$-dim Wiener increments.

\paragraph{Main theoretical results.} No CFL-type step-size restriction (the exponential
propagator handles the stiff linear part exactly). Under globally-Lipschitz $f,\sigma$:
$L^q(\Omega)$ convergence for all $q\!\ge\!2$ with temporal rate $1/4^-$ (Thm 2.3, improving
$1/8^-$ of implicit Euler--Maruyama), and almost-sure convergence (Thm 2.4). Under
non-globally-Lipschitz coefficients: convergence in probability (Thm 3.1).

\paragraph{Numerical experiments (Sec.\ 2.2.3).} $u_0(x)=\cos(\pi(x-1/2))=\sin(\pi x)$,
$f(u)=u/2$, $\sigma(u)=1-u$, $T=0.5$, $\Delta x = 2^{-9}$ ($M=512$), SEXP vs semi-implicit
Euler--Maruyama (SEM) and Crank--Nicolson--Maruyama (CNM), $\Delta t = 2^{-1}\dots 2^{-16}$,
reference $\Delta t_{\text{ref}}=2^{-16}$, $M_s=500$ samples. Fig.\ 1 shows the sup-error
$\sup_{(t,x)\in[0,0.5]\times[0,1]}\mathbb{E}[|u^{M,N}(t,x)-u^M(t,x)|^2]$ with reference-line
slope $1/2$; Fig.\ 3 illustrates almost-sure convergence.

\section{Claims table}

\begin{center}
\small
\begin{tabular}{lp{7.5cm}llll}
\toprule
ID & Claim & Type & Testable? & Tested? & Result \\
\midrule
C1 & SEXP has no CFL restriction (stable $\forall\Delta t\!\in\![2^{-1},2^{-16}]$, $M=512$)
   & numerical/stability & yes & yes & \textbf{Confirmed} \\
C2 & SEXP empirical temporal \textbf{strong order $\approx 1/2$} (Fig.\ 1)
   & numerical/quant.\ & yes & yes & \textbf{Confirmed} (RMS 0.558) \\
C3 & Almost-sure / pathwise convergence (Thm 2.4, Fig.\ 3)
   & theory + num.\ & partially & yes & \textbf{Confirmed} numerically \\
C4 & (sanity) exp integrator exact on linear part; FD $2^{\text{nd}}$-order in space
   & numerical & yes & yes & \textbf{Confirmed} (5.5e-15 / rate 2.000) \\
C5 & Improved $L^q$ rate $1/4^-$ (Thm 2.3) vs $1/8^-$ of Euler--Maruyama
   & theory & proof-only & no & Out of scope; empirical Fig-1 reproduced (C2) \\
\bottomrule
\end{tabular}
\end{center}

\section{Method}

\textbf{Environment.} CherryRd (numpy 2.4.3, scipy 1.18.0) for validation and pilots;
Monte Carlo on \texttt{uicgpu} (8$\times$A100, 255 cores; numpy 1.23.5, scipy 1.10.1) with 96
worker processes. LLM inference on the free Argo proxy at \texttt{localhost:44497}. No paid
endpoints; no \texttt{pdf} tool used.

\begin{enumerate}
  \item \textbf{Fetch OA artifact.} \texttt{curl} the arXiv PDF and LaTeX e-print;
        \texttt{pdftotext -layout}; \texttt{tar xzf} the source
        (see \texttt{artifact\_harvest.md}).
  \item \textbf{Implement SEXP from scratch} (\texttt{work/sexp\_heat.py}).
        $A = M^2 D$ is diagonalized by the discrete sine transform:
        $e^{A\Delta t}v = \mathrm{IDST}(\exp(\lambda\,\Delta t)\cdot \mathrm{DST}(v))$
        using \texttt{scipy.fft.dst/idst} type-1, \texttt{norm='ortho'}. Verified exact vs.\
        \texttt{scipy.linalg.expm}: max diff $3.3\times 10^{-16}$ at $M=16$.
  \item \textbf{Deterministic-case validation FIRST} (\texttt{validate\_deterministic.py}).
        Set $\sigma\!=\!0,\,f\!=\!0$: SEXP integrates $U'=AU$ exactly; verified
        $\Delta t$-independence ($\max|u(2^{-4})-u(2^{-10})|=\num{5.5e-15}$) and agreement
        with $e^{\lambda_1 T}\sin(\pi x)$ ($\num{1.7e-16}$). Semidiscrete FD $\to$ analytic
        PDE $u(t,x)=e^{-\pi^2 t}\sin(\pi x)$: rates $2.000, 2.000, 2.000, 2.000, 2.000$ for
        $M=32,\dots,1024$.
  \item \textbf{Strong-convergence experiment} (\texttt{run\_strong\_order\_mp.py}, uicgpu):
        exact paper parameters, coarse $\Delta t = 2^{-3},\dots,2^{-10}$,
        $\Delta t_{\text{ref}}=2^{-16}$, $M_s=500$. Coarse and fine share the SAME
        Brownian path via block-summing of the finest increments
        (Brownian-consistency), isolating temporal error.
  \item \textbf{Almost-sure/pathwise experiment} (\texttt{run\_as\_convergence.py}):
        5 independent paths, $M=512$, $\Delta t_{\text{ref}}=2^{-15}$.
  \item \textbf{Multi-judge scoring} (\texttt{run\_judges.sh}): free Argo \texttt{gpt-5.2},
        \texttt{gemini-2.5-pro}, \texttt{gpt-4.1}; opus deliberately avoided.
\end{enumerate}

\section{Results vs.\ paper}

\subsection*{C4 --- validation (analytic)}

\begin{center}
\begin{tabular}{lll}
\toprule
Check & Metric & Result \\
\midrule
SEXP linear part exact in time & $\max|u(2^{-4})-u(2^{-10})|$ & \num{5.5e-15} \\
SEXP vs $e^{\lambda_1 T}\sin(\pi x)$ & max abs & \num{1.7e-16} \\
FD $\to$ analytic, $\Delta x$ order & rate $(M=32{\to}1024)$ & 2.000, 2.000, 2.000, 2.000, 2.000 \\
\bottomrule
\end{tabular}
\end{center}

\subsection*{C1 --- stability / no CFL ($M=512$, single path, final $\max|u|$)}

\begin{center}
\small
\begin{tabular}{lccccccccc}
\toprule
$\Delta t$ & $2^{-1}$ & $2^{-2}$ & $2^{-4}$ & $2^{-6}$ & $2^{-8}$ & $2^{-10}$ & $2^{-12}$ & $2^{-14}$ & $2^{-16}$ \\
$\max|u|$ & 1.48 & 1.98 & 0.79 & 0.41 & 0.23 & 0.095 & 0.077 & 0.047 & 0.039 \\
\bottomrule
\end{tabular}
\end{center}

Bounded $\mathcal{O}(1)$ for \textbf{all} step sizes $\Rightarrow$ confirms the paper's
headline ``no CFL-type restriction.''

\subsection*{C2 --- strong temporal convergence order ($M=512$, $\Delta t_{\text{ref}}=2^{-16}$, 500 samples, $T=0.5$)}

\begin{center}
\begin{tabular}{lll}
\toprule
$\Delta t$ & $\mathbb{E}[\sup|\cdot|^2]$ & RMS $=\sqrt{\cdot}$ \\
\midrule
$2^{-3}$  & 2.575  & 1.605 \\
$2^{-4}$  & 1.121  & 1.059 \\
$2^{-5}$  & 0.487  & 0.698 \\
$2^{-6}$  & 0.218  & 0.467 \\
$2^{-7}$  & 0.101  & 0.318 \\
$2^{-8}$  & 0.048  & 0.219 \\
$2^{-9}$  & 0.0235 & 0.153 \\
$2^{-10}$ & 0.0115 & 0.107 \\
\bottomrule
\end{tabular}
\end{center}

Slope of $\mathbb{E}[\sup|\cdot|^2]$ vs.\ $\Delta t$ (log$_2$) $=1.115$; \textbf{strong (RMS)
order $=0.558$} (paper: $1/2$). Clean, consistent $1/2$-slope over 8 successive refinements.
Evidence: \texttt{evidence/evidence\_strong\_order\_full.json}.

\subsection*{C3 --- almost-sure / pathwise convergence (5 paths, $M=512$, $\Delta t_{\text{ref}}=2^{-15}$)}

Per-path sup-in-$(t,x)$ error decreases monotonically with $\Delta t$; fitted per-path slopes:
$0.568,\,0.525,\,0.562,\,0.547,\,0.535$ $\Rightarrow$ every realization converges as
$\Delta t\!\to\!0$ with order $\approx 1/2$, the numerical signature of Thm 2.4 / Fig.\ 3.
Evidence: \texttt{evidence/evidence\_as\_convergence.json}.

\section{Internal-consistency finding (noise scaling)}

A literal reading of Eq.\ \eqref{eq:sexp} --- $\Sigma=\sqrt{M}\sigma$ with
$\Delta W^n = W^M(t_{n+1})-W^M(t_n)$ and
$W^M := \sqrt{M}(W(t,x_{m+1})-W(t,x_m))\sim \mathcal{N}(0,\Delta t)$ --- applies the
$\sqrt{M}$ factor \emph{twice}, giving net per-node noise amplitude $\sqrt{M}\sigma\sqrt{\Delta t}$.
At $M=512$ this amplifies the noise by a factor $M$ and the \emph{explicit} scheme blows up for
$\Delta t > 2^{-10}$ (we measured $\max|u|\sim 10^{6}\text{--}10^{13}$), contradicting the
paper's own ``no CFL'' claim and the Fig.\ 1 x-axis range ($10^{-5}\text{--}10^{0}$).

The consistent interpretation is that the discrete increments $\Delta W^n$ are the space-time
white-noise \emph{cell} increments $\sim \mathcal{N}(0,\Delta t\cdot \Delta x)=
\mathcal{N}(0,\Delta t/M)$, so the net per-node noise is
$\sqrt{M}\sigma\sqrt{\Delta t/M}\,\xi = \sigma(u_m)\sqrt{\Delta t}\,\xi$ --- the standard,
physically-correct FD discretization of white noise, with the $\sqrt{M}$ in $\Sigma$ exactly
canceling the $1/\sqrt{M}$ in the cell increment. With this scaling the scheme is CFL-free (C1)
and reproduces Fig.\ 1 (C2). This is a presentation ambiguity in the paper, not a defect in the
method; all three LLM judges considered the reconciliation acceptable and standard.

\section{Judge assessments (free Argo, non-opus)}

\begin{center}
\begin{tabular}{ll}
\toprule
Judge & Verdict \\
\midrule
\texttt{argo:gpt-5.2}          & \textbf{REPLICATED} (``\ldots strong-rate estimate consistent with $\approx 1/2$; reconciliation acceptable'') \\
\texttt{argo:gemini-2.5-pro}   & \textbf{REPLICATED} (``\ldots 0.558 a convincing match; reconciliation is a strength'') \\
\texttt{argo:gpt-4.1}          & \textbf{REPLICATED} (``\ldots matches all headline claims; noise-scaling reconciliation is not a concern'') \\
\bottomrule
\end{tabular}
\end{center}

\section{Files}

\begin{itemize}
  \item \texttt{work/sexp\_heat.py} --- SEXP solver (DST-diagonalized exponential integrator).
  \item \texttt{work/validate\_deterministic.py} --- analytic-case validation.
  \item \texttt{work/run\_strong\_order.py}, \texttt{work/run\_strong\_order\_mp.py} --- strong-order (serial / MP).
  \item \texttt{work/run\_as\_convergence.py} --- almost-sure/pathwise experiment.
  \item \texttt{work/run\_judges.sh}, \texttt{work/judge\_summary.txt} --- multi-judge harness + prompt.
  \item \texttt{report/evidence/} --- JSON/text outputs, judge responses, validation logs.
  \item \texttt{work/anton\_cohen\_2015.pdf}, \texttt{work/src/} --- OA paper + LaTeX source.
\end{itemize}

\section{Assessment}

Every reproducible headline claim was independently confirmed with a from-scratch
implementation: the exponential integrator is exact for the linear part (validated to
machine precision), the spatial FD operator is exactly 2nd order, the scheme is CFL-free,
and the temporal \textbf{strong convergence order matches the paper's $1/2$} (measured
$0.558$ over 500 samples at the paper's exact parameters), with pathwise convergence also
confirmed. The only theoretical claim not reproduced is the proof-level $1/4^-$ $L^q$ rate
(Thm 2.3), which is analytical rather than a benchmarkable number; the corresponding
\emph{empirical} Fig-1 order ($1/2$) is reproduced. One presentation-level noise-scaling
ambiguity was identified and reconciled without affecting the conclusions.

% ==========================================================================
\section{GENUINE CRITIQUE}
% ==========================================================================

This section deliberately steps outside the ``verdict-defending'' mode of the report body and
records the strongest good-faith objections a skeptical reviewer could raise against our
replication, together with our honest assessment of how much each one bites. The verdict
remains \textsc{Replicated}, but readers should weigh these caveats before treating this
report as a certification of the paper.

\subsection*{1. Empirical order vs.\ proof-level order}
We match the \emph{empirical Fig.\ 1 rate} of $1/2$ (measured $0.558$), not the paper's
\emph{proved} $L^q$ rate of $1/4^-$ (Thm 2.3). These are genuinely different objects: the proof
uses an $L^q(\Omega)$-in-$C([0,T]\times[0,1])$ norm and includes a spatial discretization
contribution; Fig.\ 1 (and our C2) fixes the space grid at $M=512$ and measures only the
\emph{temporal} error using a fine reference on the same grid. The observed $1/2$ is a
well-known ``super-convergence''-type effect for temporal-only refinement of this class of
schemes and does not contradict the proved rate --- but a reader who conflates ``we match
Fig.\ 1'' with ``we replicate Thm 2.3'' would be over-claiming, and we have not attempted to
reproduce the proof-level rate under joint $(\Delta t,\Delta x)$ refinement.

\subsection*{2. Noise-scaling reconciliation is an interpretation, not a proof}
Section 5 documents that a literal-symbol reading of Eq.\ (15) blows up, and that the
standard cell-increment interpretation ($\Delta W^n\sim \mathcal{N}(0,\Delta t/M)$) restores
everything. Three judges accepted this as ``the obvious reading.'' However, we could not point
to a single sentence in the paper that unambiguously specifies which convention is meant; the
reconciliation is inferred from consistency with the paper's other claims (CFL-freeness,
Fig.\ 1 range). A less charitable reviewer could argue that we ``chose the convention that
makes the paper right,'' i.e.\ that our replication is partly retrofitted. We think this is
mitigated by the fact that the chosen convention is the physically standard one and is what any
practitioner in SPDE numerics would implement --- but it is a real methodological caveat.

\subsection*{3. Single test problem}
Both the paper's Fig.\ 1 and our C2 use exactly one triple $(u_0,f,\sigma) =
(\sin(\pi x),\,u/2,\,1-u)$. We did not vary these to check robustness of the observed $1/2$
rate. In particular $\sigma(u)=1-u$ is affine and globally Lipschitz; a stress test with
nonlinear or nearly-non-Lipschitz $\sigma$ (as in Sec.\ 3 of the paper) would be a
substantially stronger replication and is genuinely absent.

\subsection*{4. Single $M$ and single $T$}
The strong-order experiment is at $M=512$, $T=0.5$ only. We did not check that the measured
$0.558$ slope is invariant to $M$ (e.g.\ $M\in\{128,256,1024\}$) or to $T$ (e.g.\ $T\in\{1,2\}$).
The paper's Fig.\ 1 uses the same $M$ and $T$, so this is a fair replication of what was
plotted, but it is not a stress test of the rate.

\subsection*{5. Sample size and Monte Carlo error}
$M_s=500$ matches the paper. The empirical slope $0.558$ vs.\ theoretical $0.500$ leaves
$\approx 0.06$ of ``slack''; we did not compute a bootstrap CI on the slope, and it is
plausible that the true slope is exactly $1/2$ with the observed excess being MC noise
(especially at the coarsest $\Delta t=2^{-3}$ where $\mathbb{E}[\sup|\cdot|^2]=2.575$ is the
largest and least stable estimator). A cleaner replication would report a slope $\pm$ 95\% CI
rather than a point estimate.

\subsection*{6. Two-scheme comparison omitted}
Fig.\ 1 in the paper is a \emph{comparison} of SEXP against SEM and CNM. We implemented and
measured only SEXP. The relative claim ``SEXP is competitive with / better than SEM, CNM'' is
therefore \emph{not} independently verified by us --- we only checked the SEXP curve. A full
replication of Fig.\ 1 would reimplement all three and reproduce the ranking.

\subsection*{7. Brownian-consistency by block-summing is correct but subtle}
Our coarse/fine paths share the same Brownian path via block-summing of the finest increments.
This is standard and correct, but it is also the single most bug-prone piece of any SPDE
strong-error experiment. We validated it by construction (the finest coarse level with
$\Delta t=2^{-10}$ still recovers order $\approx 1/2$ against the $\Delta t=2^{-16}$ reference)
but did not do a paranoid cross-check against an independent PRNG-seeded pairing.

\subsection*{8. Judges see only our summary}
The three LLM judges saw the quantitative summary we wrote, not the underlying JSON or code.
Their agreement therefore reflects internal coherence of \emph{our narrative}, not independent
audit of \emph{our code}. Their votes should be read as ``the write-up is self-consistent'' not
as ``the implementation is correct.'' The strongest evidence remains the deterministic
validation (5.5e-15 time-exactness, rate 2.000 in space).

\subsection*{9. Verdict robustness}
Even after applying every caveat above, the deterministic validation is machine-precision, the
strong-order slope is monotone and clean over 8 refinements, the pathwise slopes are tightly
clustered around $0.5$, and the qualitative Fig.\ 1 story reproduces cleanly. We therefore
keep the verdict at \textsc{Replicated}, but a fair reading of ``replicated'' here is:
\emph{``the paper's SEXP method, at the paper's own test problem and parameters, and under
the standard cell-increment interpretation of the discrete noise, reproduces the
empirical Fig.\ 1 order-$1/2$ rate and the CFL-freeness claim, from an independent
from-scratch implementation, with three free-endpoint LLM judges concurring.''}
It is \emph{not}: ``we independently proved Thm 2.3'' or ``we ranked SEXP against SEM/CNM.''

\section*{Verdict}
\textbf{REPLICATED.}

\end{document}
